\begin{table}[htbp]
  \centering
  \caption{Decomposition of Sectoral Labor Productivity Growth, 1950--2023}
  \label{tab:ex1_decomp}
  \begin{threeparttable}
  \centering
  \begin{tabular}{l *{3}{S[table-format=-1.2]}}
    \toprule
     & \multicolumn{3}{c}{Sector} \\
    \cmidrule(lr){2-4}
     & {Goods} & {RI-Serv.} & {NRI-Serv.} \\
    \midrule
    I.\ Endogenous innovation & 2.05 & 0.82 & 0.71 \\
    II.\ Embodied technical progress & 0.35 & 1.00 & 0.45 \\
    III.\ Relative price effect & -0.15 & 0.13 & 0.14 \\
    IV.\ Labor efficiency & 0.37 & 0.34 & 0.34 \\
    \midrule
    Total & 2.62 & 2.29 & 1.64 \\
    \bottomrule
  \end{tabular}
  \begin{tablenotes}[flushleft]
    \footnotesize
    \RaggedRight
    \item[] \textit{Notes:} Entries report annualized contributions, in percentage points, to the growth of sectoral labor productivity $\mathcal{A}_j = Y_j/L_j^p$ over 1950--2023. The channels follow Equation (18): endogenous innovation is $e_j^Y(t) g_{N_j}(t)$, embodied technical progress reflects machine productivity, the relative price effect reflects sectoral output prices. Entries are rounded so that they add to the printed totals.
  \end{tablenotes}
  \end{threeparttable}
\end{table}
